\documentclass[journal=jpclcd,manuscript=letter]{achemso}

\usepackage[version=4]{mhchem}
\usepackage{siunitx}
\usepackage{booktabs}
\usepackage{array}
\usepackage{graphicx}
\usepackage{caption}
\usepackage{subcaption}
\usepackage{amsmath}
\usepackage{amssymb}
\usepackage{orcidlink}
\hypersetup{hidelinks}
\usepackage{pgfplots}
\pgfplotsset{compat=1.18}
\usepackage{tikz}
\usetikzlibrary{patterns,calc}
\usetikzlibrary{arrows.meta,positioning,fit,backgrounds,shapes.geometric}
\usepackage{float}
\usepackage{placeins}
\raggedbottom
\setlength{\textfloatsep}{8pt plus 2pt minus 2pt}
\setlength{\floatsep}{6pt plus 2pt minus 2pt}
\setlength{\intextsep}{8pt plus 2pt minus 2pt}
\setlength{\abovecaptionskip}{4pt}
\setlength{\belowcaptionskip}{2pt}
\hfuzz=20pt
\vfuzz=5pt
\vbadness=10000
\setlength{\emergencystretch}{2em}
\makeatletter
\setlength{\acs@tocentry@width}{4.35in}
\setlength{\acs@tocentry@height}{2.35in}
\makeatother

\newcommand{\angstrom}{\ensuremath{\text{\AA}}}
\newcommand{\morpheus}{\textsc{Morpheus}}

\author{Vincenzo Barone\,\orcidlink{0000-0001-6420-4107}}
\affiliation{Consorzio Interuniversitario Nazionale per la Scienza e Tecnologia dei Materiali (INSTM), via G. Giusti 9, 50121 Firenze, Italy}
\email{vincebarone52@gmail.com}

\title{Transferable Class-Correction Refinement from Multiple Rotational Data Sets}

\begin{document}
\raggedbottom

\begin{tocentry}
\centering
\resizebox{0.94\linewidth}{!}{\begin{tikzpicture}[x=1cm,y=1cm,>=Latex]
\tikzstyle{box}=[draw,rounded corners=3pt,align=center,minimum height=0.75cm]
\node[box,fill=blue!8,minimum width=2.6cm] (m1) at (0,2.4) {Molecule 1\\topology, symmetry};
\node[box,fill=blue!8,minimum width=2.6cm] (m2) at (0,1.2) {Molecule 2\\topology, symmetry};
\node[box,fill=blue!8,minimum width=2.6cm] (m3) at (0,0.0) {Molecule $n$\\topology, symmetry};
\node[box,fill=green!10,minimum width=2.9cm] (coords) at (3.4,1.2) {non-redundant\\GIC spaces};
\node[box,fill=yellow!16,minimum width=2.9cm] (classes) at (6.8,1.2) {shared chemical\\class map $T_{m,kc}$};
\node[box,fill=orange!15,minimum width=2.6cm] (fit) at (10.0,1.2) {global fit\\$\Delta_c$};
\draw[->] (m1.east) -- (coords.west);
\draw[->] (m2.east) -- (coords.west);
\draw[->] (m3.east) -- (coords.west);
\draw[->] (coords.east) -- (classes.west);
\draw[->] (classes.east) -- (fit.west);
\node[align=center,font=\small] at (6.8,0.1) {soft priors enter as\\Gaussian observations};
\node[align=center,font=\small] at (10.0,0.1) {rank, covariance,\\condition number};
\end{tikzpicture}
}
\end{tocentry}

\begin{abstract}
We introduce a multi-molecule semiexperimental refinement strategy in which rotational data from chemically related molecules are combined through transferable coordinate-class corrections.  Each molecule keeps its own topology, symmetry, isotopic substitutions, vibrational corrections, and non-redundant coordinate basis, while the fitted variables are corrections shared by chemically defined classes such as carbonyl C--O, single C--O, C--C, C--H, or selected angular families.  Classes may combine atom symbols, primitive-coordinate values, and continuous synthon-derived atom types.  The approach replaces molecule-specific model construction by a rank-checked global statistical model that can use sparse isotopic information across homologous series.  Maleic, phthalic, and succinic anhydrides provide a chemically controlled homologous-series test, while the two glycine conformers test whether the same class definitions can transfer across different molecular geometries and atom orderings.  The comparison of no-prior, soft-prior, and hard-constraint variants shows that formally admissible angular classes can be severely ill conditioned without statistical regularization, whereas Gaussian soft priors retain the chemically explicit model while restoring numerical stability.
\end{abstract}

\section*{Introduction}

Semiexperimental structure refinement is usually formulated molecule by molecule.  This is natural when isotopic substitution is rich enough to determine an independent structure, but it becomes inefficient for larger molecules or homologous series where different members share the same systematic quantum-chemical errors.  In those cases the transferable information is not an absolute bond length or angle, but a correction relative to the computational reference within a chemically defined coordinate class.

Here we formulate this idea as a global refinement problem.  For molecule $m$ and generated internal coordinate $q_k$,
\begin{equation}
q_{m,k}^{\mathrm{SE}} =
q_{m,k}^{\mathrm{QC}} + T_{m,kc}\Delta_c ,
\end{equation}
where $T_{m,kc}$ maps the generated coordinate onto class $c$ and $\Delta_c$ is shared across all molecules.  The mapping is coefficient based: if a symmetry-adapted coordinate contains several primitive stretches, bends, torsions, or out-of-plane terms, the class correction is projected through the primitive coefficients rather than assigned by a binary label.

\begin{figure}[H]
\centering
\resizebox{0.96\linewidth}{!}{\begin{tikzpicture}[x=1cm,y=1cm,>=Latex]
\tikzstyle{box}=[draw,rounded corners=3pt,align=center,minimum height=0.75cm]
\node[box,fill=blue!8,minimum width=2.6cm] (m1) at (0,2.4) {Molecule 1\\topology, symmetry};
\node[box,fill=blue!8,minimum width=2.6cm] (m2) at (0,1.2) {Molecule 2\\topology, symmetry};
\node[box,fill=blue!8,minimum width=2.6cm] (m3) at (0,0.0) {Molecule $n$\\topology, symmetry};
\node[box,fill=green!10,minimum width=2.9cm] (coords) at (3.4,1.2) {non-redundant\\GIC spaces};
\node[box,fill=yellow!16,minimum width=2.9cm] (classes) at (6.8,1.2) {shared chemical\\class map $T_{m,kc}$};
\node[box,fill=orange!15,minimum width=2.6cm] (fit) at (10.0,1.2) {global fit\\$\Delta_c$};
\draw[->] (m1.east) -- (coords.west);
\draw[->] (m2.east) -- (coords.west);
\draw[->] (m3.east) -- (coords.west);
\draw[->] (coords.east) -- (classes.west);
\draw[->] (classes.east) -- (fit.west);
\node[align=center,font=\small] at (6.8,0.1) {soft priors enter as\\Gaussian observations};
\node[align=center,font=\small] at (10.0,0.1) {rank, covariance,\\condition number};
\end{tikzpicture}
}
\caption{Multi-molecule class-correction refinement.  Each molecule contributes its own topology-derived non-redundant coordinate space; chemically defined class maps connect these local spaces to a shared set of global corrections.}
\label{fig:model_scheme}
\end{figure}

\section*{Model}

The coordinate classes are type aware.  Stretches use unordered atom pairs and may be refined by value windows, which separates, for example, carbonyl and single C--O corrections without introducing molecule-specific rules.  Bends preserve the central atom.  Torsions are classified by the two central atoms, i.e., by the central bond.  Out-of-plane coordinates are classified by their central atom.  When element symbols are too coarse, atomic synthons provide a continuous local-environment descriptor.  In the present implementation this information is condensed into an effective atomic number, $Z_\mathrm{eff}$, and two atoms are assigned to the same class when their $Z_\mathrm{eff}$ values differ by less than a chosen threshold.  This prevents chemically different relations from being merged by a purely textual or fully sorted atom list.

Soft priors are ordinary Gaussian observations on the class correction,
\begin{equation}
\Delta_c = \Delta_c^0 \pm \sigma_c ,
\end{equation}
and enter the same weighted least-squares problem as the spectroscopic rows.  They are therefore weaker than hard constraints but stronger than unregularized weak directions.  The distinction is essential for angular classes in isotope-poor data sets.
The global linear problem is solved in a column-scaled, rank-revealing singular-value basis; the same truncation threshold is used for the correction vector, covariance matrix, numerical rank, condition number, and reported correlation diagnostics.
Each fit is then assigned an explicit model status.  Models are rejected when they are rank deficient, singular, severely ill conditioned, unsupported by the requested molecule coverage, or dominated by nearly identical class directions.  Less severe but chemically informative correlations are reported as review items rather than hidden in the covariance matrix.

\section*{Anhydride Test Set}

Maleic, phthalic, and succinic anhydrides were selected as a homologous set with literature rotational data and different isotopic completeness.  The same class model was applied to all three molecules.  Stretch classes were left unregularized, whereas broad bend classes were tested in three variants: no prior, soft prior, and hard constraint.

\begin{table}[H]
\centering
\caption{Coordinate classes used for the anhydride ensemble refinement.  Stretch classes distinguish short/long C--C and carbonyl/single C--O ranges.  Bend priors are Gaussian soft observations on the class correction, not hard frozen coordinates.}
\label{tab:anhydride_classes}
\begin{tabular}{llll}
\toprule
Class & Coordinate kind & Selector & Prior \\
\midrule
CC\_short & stretch & C--C, $r\leq1.42$ & none \\
CC\_long & stretch & C--C, $r\geq1.42$ & none \\
CO\_carbonyl & stretch & C--O, $r\leq1.28$ & none \\
CO\_single & stretch & C--O, $r\geq1.28$ & none \\
CH\_stretch & stretch & C--H & none \\
CCC\_bend & bend & C--C--C & $0.0 \pm 1.00\times 10^{-3}$ \\
CCO\_bend & bend & C--C--O & $0.0 \pm 1.00\times 10^{-3}$ \\
COC\_bend & bend & C--O--C & $0.0 \pm 1.00\times 10^{-3}$ \\
OCO\_bend & bend & O--C--O & $0.0 \pm 1.00\times 10^{-3}$ \\
\bottomrule
\end{tabular}

\end{table}

\begin{table}[H]
\centering
\caption{Class-support diagnostics for the soft-prior anhydride model.  The prior residual ratio measures how far the fitted correction moves from the prior centre in units of the prior standard deviation.}
\label{tab:class_support}
\begin{tabular}{llrrr}
\toprule
Class & Kind & Matched coordinates & Molecules & $|\mathrm{prior\ residual}|/\sigma_{\mathrm{prior}}$ \\
\midrule
CC\_short & stretch & 14 & 2 & -- \\
CC\_long & stretch & 15 & 3 & -- \\
CO\_carbonyl & stretch & 10 & 3 & -- \\
CO\_single & stretch & 4 & 3 & -- \\
CH\_stretch & stretch & 7 & 3 & -- \\
CCC\_bend & bend & 12 & 3 & 0.002 \\
CCO\_bend & bend & 12 & 3 & 0.011 \\
COC\_bend & bend & 12 & 3 & 0.003 \\
OCO\_bend & bend & 7 & 3 & 0.005 \\
\bottomrule
\end{tabular}

\end{table}

\begin{figure}[H]
\centering
\begin{tikzpicture}
\begin{axis}[
  width=0.82\linewidth,
  height=5.2cm,
  xmode=log,
  xlabel={Scaled condition number},
  ylabel={Weighted RMS after fit},
  grid=both,
  legend style={draw=none,fill=none,at={(0.03,0.03)},anchor=south west},
]
\addplot+[only marks,mark=*,mark size=2.6pt] coordinates {(428628992.313,0.00760258786126)};
\addlegendentry{No Prior}
\addplot+[only marks,mark=square*,mark size=2.6pt] coordinates {(69.887278089,0.110941272264)};
\addlegendentry{Soft Prior}
\addplot+[only marks,mark=triangle*,mark size=2.6pt] coordinates {(69.4860339176,0.111107899864)};
\addlegendentry{Hard Constraint}
\end{axis}
\end{tikzpicture}

\caption{Numerical stability and residual reduction for the three anhydride variants.  The unregularized angular model reaches a lower residual but has a catastrophic scaled condition number; the soft-prior model retains the angular classes while recovering the numerical stability of the hard-constrained model.}
\label{fig:condition_wrms}
\end{figure}

\begin{figure}[H]
\centering
\begin{subfigure}{0.49\linewidth}
\centering
\begin{tikzpicture}
\begin{axis}[
  width=0.82\linewidth,
  height=5.2cm,
  xmode=log,
  xlabel={Bend prior $\sigma$},
  ylabel={Weighted RMS after fit},
  grid=both,
]
\addplot+[mark=*,mark size=2pt] coordinates {
  (1e-05,0.0677562551194)
  (3e-05,0.0677562129428)
  (0.0001,0.0677557331875)
  (0.0003,0.0677515158819)
  (0.001,0.0677035848336)
  (0.003,0.067285417253)
  (0.01,0.0629025553883)
  (0.03,0.0412523294019)
  (0.1,0.0180612261556)
};
\end{axis}
\end{tikzpicture}

\caption{Residual response}
\end{subfigure}
\begin{subfigure}{0.49\linewidth}
\centering
\begin{tikzpicture}
\begin{axis}[
  width=0.82\linewidth,
  height=5.2cm,
  xmode=log,
  ymode=log,
  xlabel={Bend prior $\sigma$},
  ylabel={Scaled condition number},
  grid=both,
]
\addplot+[mark=*,mark size=2pt] coordinates {
  (1e-05,79.0830479026)
  (3e-05,79.0832956179)
  (0.0001,79.0861117602)
  (0.0003,79.1107420515)
  (0.001,79.3762402621)
  (0.003,81.0644975565)
  (0.01,88.4527308818)
  (0.03,102.518121878)
  (0.1,127.784574141)
};
\end{axis}
\end{tikzpicture}

\caption{Conditioning response}
\end{subfigure}
\caption{Prior-strength scan for the angular classes.  Increasing the prior width gradually lowers the residual but also worsens conditioning; the default value $\sigma=10^{-3}$ lies near the stable end of this tradeoff.}
\label{fig:prior_scan}
\end{figure}

\begin{table}[H]
\centering
\caption{Systematic comparison of ensemble class-correction models for the three anhydrides.  Hard constraints correspond to removing the weak angular corrections from the global variable set.  The status column is assigned by the rank/conditioning/support policy described in the text.}
\label{tab:anhydride_prior_comparison}
\begin{tabular}{llrrrrr}
\toprule
Variant & Status & Classes & Rank & Cond. & WRMS$_0$ & WRMS \\
\midrule
No Prior & rejected & 8 & 8 & $4.29\times 10^{8}$ & 0.18437 & 0.0076 \\
Soft Prior & review & 8 & 8 & 69.88728 & 0.18437 & 0.11094 \\
Hard Constraint & review & 4 & 4 & 69.48603 & 0.18437 & 0.11111 \\
\bottomrule
\end{tabular}

\end{table}

The no-prior model has the lowest residual but is rejected by the acceptance policy because the design matrix is ill conditioned and contains nearly indistinguishable angular directions.  The soft-prior model restores the condition number to the same scale as the hard-constrained model while retaining explicit angular classes and reporting their posterior residuals.  It remains in review because the carbonyl and single C--O corrections are strongly anticorrelated, a useful chemical diagnostic rather than a hidden failure.  Stability and predictivity are therefore kept as separate diagnostics: condition number and posterior prior residuals test whether the model is numerically admissible, whereas leave-one-molecule-out residuals test whether a correction learned on two molecules transfers to the third.

\begin{table}[H]
\centering
\caption{Leave-one-molecule-out prediction with the soft-prior model.  The transfer score is the predicted WRMS divided by the uncorrected held-out WRMS.  The result is diagnostic rather than conclusive: the three-member anhydride set is sufficient to expose numerical regularization, but too small to establish broad predictive transferability.}
\label{tab:loo}
\begin{tabular}{lrrrrr}
\toprule
Held-out & Train & Cond. & WRMS$_0$ & WRMS$_\mathrm{pred}$ & Score \\
\midrule
Maleic Anhydride & 2 & 165.49 & 0.0031 & 0.14311 & 46.21 \\
Phthalic Anhydride & 2 & 127.18 & 0.15162 & 2.75495 & 18.17 \\
Succinic Anhydride & 2 & 120.88 & 0.34259 & 0.37161 & 1.08 \\
\bottomrule
\end{tabular}

\end{table}

The leave-one-molecule-out test is deliberately severe.  Maleic and succinic anhydrides remain qualitatively stable when predicted from the other two members, whereas phthalic anhydride is not predicted well from maleic and succinic alone.  This failure mode is useful: it shows that the ensemble model reports lack of transferability rather than hiding it, and motivates extension to larger homologous families before predictive claims are made.

\section*{Glycine Conformer Test}

The anhydride comparison tests transfer across a homologous chemical family.  A complementary question is whether class definitions can be transferred between different conformers of the same molecule, where the chemical graph is conserved but the geometry, symmetry, atom ordering, and isotopic information are not identical.  Glycine conformers I and II provide such a case.  The two calculations used the existing semiexperimental benchmark inputs directly; the ensemble layer builds an independent non-redundant GIC model for each conformer and then maps both models to the same global class corrections.

\begin{table}[H]
\centering
\caption{Glycine I/II ensemble refinement using two class definitions.  The element/value model uses ordinary atom symbols and C--O value windows.  The synthon model replaces those coarse atom types by continuous $Z_\mathrm{eff}$ prototypes with a threshold of 0.035.}
\label{tab:glycine_typing}
\begin{tabular}{lrrrr}
\toprule
Class definition & Rank & Scaled condition & WRMS before & WRMS after \\
\midrule
Element and value windows & 11 & 721.08 & 4900.88 & 29.94 \\
Continuous synthon $Z_\mathrm{eff}$ types & 11 & 710.11 & 4900.88 & 24.90 \\
\bottomrule
\end{tabular}

\end{table}

Both models are rank complete, but the synthon-typed model lowers the weighted RMS from 29.94 to 24.90 while keeping essentially the same scaled condition number.  The improvement is modest but important mechanistically: it is obtained without manually identifying corresponding atoms in the two conformers.  The carboxyl carbon, methylene carbon, carbonyl oxygen, hydroxyl oxygen, amino hydrogens, and methylene hydrogens are recognized by their continuous local environments, not by atom labels.

\begin{table}[H]
\centering
\caption{Threshold scan for the continuous glycine synthon atom types.  The model is accepted and unchanged from 0.010 to 0.075, showing that the selected threshold is inside a stable plateau.  At 0.100, chemically distinct atom types begin to overlap, the numerical rank drops, and the acceptance policy rejects the singular model.}
\label{tab:glycine_synthon_threshold}
\begin{tabular}{rlrrrr}
\toprule
$Z_\mathrm{eff}$ threshold & Status & Rank & Cond. & WRMS & Min. matches \\
\midrule
0.010 & accepted & 10 & 227.15 & 25.96 & 3 \\
0.035 & accepted & 10 & 227.15 & 25.96 & 3 \\
0.075 & accepted & 10 & 227.15 & 25.96 & 3 \\
0.100 & rejected & 9 & $\infty$ & 33.09 & 3 \\
0.150 & rejected & 9 & $\infty$ & 33.09 & 3 \\
0.250 & rejected & 9 & $\infty$ & 33.09 & 3 \\
\bottomrule
\end{tabular}

\end{table}

\begin{table}[H]
\centering
\caption{Continuous synthon classes used for the glycine I/II fit.  The prototype column lists the $Z_\mathrm{eff}$ values defining the class; each primitive coordinate matches when the chemically relevant $Z_\mathrm{eff}$ signature is within 0.035 of the prototype.  Corrections are in internal-coordinate units.}
\label{tab:glycine_synthon_classes}
\begin{tabular}{llrrr}
\toprule
Class & Type & $Z_\mathrm{eff}$ prototype & Matches & Correction \\
\midrule
Ccarb--Cmethylene & stretch & 5.91, 5.88 & 8 & -0.00507 \\
Cmethylene--N & stretch & 5.88, 6.96 & 8 & -0.00220 \\
Ccarb--Ocarbonyl & stretch & 5.91, 7.865 & 8 & -0.00250 \\
Ccarb--Ohydroxyl & stretch & 5.91, 7.95 & 9 & -0.00542 \\
Cmethylene--H & stretch & 5.88, 0.838 & 6 & -0.00280 \\
N--H & stretch & 6.96, 0.868 & 5 & -0.00323 \\
Ohydroxyl--H & stretch & 7.95, 0.898 & 9 & 0.00735 \\
Ccarb--Cmethylene--Ocarbonyl & bend & 5.88, 5.91, 7.865 & 16 & -0.00553 \\
Ccarb--Cmethylene--N & bend & 5.91, 5.88, 6.96 & 3 & -0.00223 \\
H--Cmethylene--H & bend & 0.838, 5.88, 0.838 & 13 & -0.00301 \\
H--N--H & bend & 0.868, 6.96, 0.868 & 16 & -0.01906 \\
\bottomrule
\end{tabular}

\end{table}

The glycine result therefore tests a different part of the formalism from the anhydrides.  The anhydrides show that soft priors can regularize weak class directions in a sparse homologous set.  Glycine shows that transferable classes do not have to be hand-coded as molecule-specific atom lists: continuous synthon types provide a route to chemically meaningful class definitions across conformers.

\section*{Outlook}

The present letter skeleton establishes two complementary proof-of-principle tests: a homologous anhydride series and a conformer-pair glycine test.\cite{morpheus_ensemble_working}  The next validation layer should add a larger homologous family, for example steroid hormones, where sparse isotopic data can be combined with shared class corrections.  The same formalism is also relevant beyond semiexperimental structures, including force-field calibration, systematic correction of quantum-chemical geometries, and transferability tests across molecular families.

\bibliography{references}

\end{document}
